% Standalone local theorem and proof; requires amsmath,amssymb,amsthm.
% This is a conditional local exclusion, not a replacement for the structural
% theorem or for the actual activity/retention classification.
\subsection{An analytic exclusion of the first-gap-one A0 cross case}

Consider the normalized cardinality-$85$ source with $d=77$ and
$E=P_s\cup\{0\}\cup(-P_t)$, where $r=|P_s|$, $k=|P_t|$ and
$r+k=68$. Assume that the first missing positive subset sum of the
absolute core is $1$, that $z=w=2$, and that the A0 classification gives
the retained centers $P=C_{4s}$, $\alpha_t$, $C_{6t}$ and $U_s$ and all
eight CORE8 resources. The selected $P$ is an actual contributor of its
sole-owned paired unit. The fixed-label credit is $9$; write
$\varepsilon=\delta-9\ge0$.

The inputs used here are the proved source bound
\[
 T=26+I+\tau+\delta\le165,\qquad I=(k-34)^2,
\]
the nominal first-gap-one endpoint bounds
\[
 Z\ge r_o+4-\ell_Z,\qquad Z\le S_v-r_v-3+h_Z,
\]
and the actual native/grouped loss budgets. All subset copies used below
belong to the specified retained resources; no capacity is assigned to
an empty center.

Write the four increasing source-tail offsets as $a_v<b_v<c_v<e_v$.
Let $U=a_s+b_s+a_t+b_t$, $J=e_s-c_s+e_t-c_t$, and let $B_v$ count
the distinct sums of two offsets on sign $v$ that exceed $e_v$.
Set $B=B_s+B_t$ and let $p$ be the number of odd integers
$a_v+b_v$. The two source envelopes and the quartet fringe bound give
\[
 J+U\le332-T,\quad J\ge T-18+p,\quad
 2B\le U+16-p.
\]
Thus $2T+U+p\le350$ and
\begin{equation}\label{eq:85-cross-fringe}
 B\le\min(20,183-T).
\end{equation}
The full-source central interval gives
$K\le M+R+1+B$, where $M=M_s+M_t$ and $R=R_s+R_t$.
Consequently the exact original mass identity yields
\begin{equation}\label{eq:85-cross-source}
 S_s+S_t=2K+1164-\delta
 \le2M+2R+1166+2B-\delta.
\end{equation}

\paragraph{The endpoint transport, including every constant.}
Put $a=s_4=M_s+X_s$, $b=t_4=M_t+X_t$ and $G=g_s+g_t$.
The literal center identity is
\[
 U_s=P+\alpha_t+C_{6t}+2a+4b+2G.
\]
Three lower endpoints and the upper endpoint of $U_s$ imply
\[
 S_s\ge219-2k+2a+4b+2G-E_{\rm end},\qquad
 E_{\rm end}=\ell_P+\ell_{\alpha_t}+\ell_{C_{6t}}+h_{U_s}
 \le\varepsilon.
\]
Indeed their constant is
$(k+4)+2(r+4)+(r+3)=219-2k$.
The opposite source satisfies
$S_t\ge T_k+\tau_t+4R_t+10$, while
$R_t\ge k+1$ and $M_t\ge k+5$.
Combining with \eqref{eq:85-cross-source} therefore gives
\begin{equation}\label{eq:85-cross-before-overlap}
 T_k+4k+\tau_t+243+2X_s+4X_t+2G
 \le2R+1166+2B-\delta+E_{\rm end}.
\end{equation}
Here $243=219+10+14$, and $14$ comes from
$2M_t+4R_t-2k\ge4k+14$.

\paragraph{Admissible asymmetric core copies.}
Put
\[
 E_v=\tau_v-R_v+1\ge0,\quad
 \zeta=\tau_t-k\ge0,\quad D=E_s+E_t=\tau-R+2.
\]
Every core magnitude is at least its rank plus $1$.
After deleting the maximum, its remaining rank deviations sum to $E_v$.
Deleting the $j_v$ largest core entries therefore leaves maximum at most
$r_v-j_v+1+\lfloor E_v/j_v\rfloor$.
For the signed trimmed core $Q$, this proves
\[
 |Q\cap(Q+q)|\ge
 67-j_s-j_t-\lfloor E_s/j_s\rfloor-\lfloor E_t/j_t\rfloor-q.
\]
Choose $j_s=3$ and choose $j_t\in\{1,2,3\}$ minimizing
$j_t+\lfloor E_t/j_t\rfloor$. Define
\[
 \phi_s=3+\lfloor E_s/3\rfloor,\quad
 \phi_t=\min_{j=1,2,3}\bigl(j+\lfloor E_t/j\rfloor\bigr),\quad
 f=67-\phi_s-\phi_t,
 \quad F=\max(137-R,f).
\]
The full core gives the bound $137-R-q$; the trimmed core gives $f-q$.

For a leading-to-own3 edge, the leading resource contains
$\lambda_v+M_v+Q$ and the ordinary resource contains
$\lambda_v+v_4+Q$, with gap $X_v=v_4-M_v$.
The reflected leading-to-$C_7$ edge uses the reflected copies at the
same gap. Thus these are actual subsets of the retained resources.
The rooted forest count, separately for an individual two-edge star
and for the four-edge star, gives
\[
 2X_t\ge2F-\varepsilon,\qquad
 2X_s+2X_t\ge4F-\varepsilon,
\]
and hence $2X_s+4X_t\ge6F-2\varepsilon$.
The four disjoint CORE8 ordinary pairs similarly give
$2G\ge4f-\varepsilon$.
Substituting into \eqref{eq:85-cross-before-overlap} and using
$E_{\rm end}\le\varepsilon$ and $\varepsilon=\delta-9$ yields
\begin{equation}\label{eq:85-cross-raw}
 T_k+4k+\tau_t+6F+4f
 \le2R+887+2B+3\delta.
\end{equation}
In particular the conservative old-$B=20$ residual is
\[
 L=T_k+4k+\tau_t-927-2R-3\delta+6F+4f.
\]
The actual residual in \eqref{eq:85-cross-raw} is
$L+2(20-B)$; this is exactly the fringe correction, not an additional
credit against $\delta$.

\paragraph{A short integer certificate.}
Put $q=T_k+5k$. Since $T=26+I+\tau+\delta$,
\[
 L+3T=q+3I-855+\zeta+3(E_s+E_t)+R+6F+4f.
\]
The two inequalities $F\ge137-R$ and $F\ge f$ give
$6F\ge137-R+5f$, so
\[
 L+3T\ge q+3I-115
 +(3E_s-9\phi_s)+(\zeta+3E_t-9\phi_t).
\]
First, $3E_s-9\phi_s\ge-27$.
If $E_t=0$, the second parenthesis is at least $-9$.
If $E_t>0$, then $R_t-k-1\ge1$: equality would force the entire
$k$-element core to be $\{2,\ldots,k+1\}$, for which $E_t=0$.
Hence $\zeta=E_t+(R_t-k-1)\ge E_t+1$.
For $E_t\ge6$, the choice $j_t=3$ gives
\[
 \zeta+3E_t-9\phi_t\ge E_t-26\ge-20.
\]
For $E_t=1,2,3,4,5$, the respective bounds are
$-13,-18,-14,-19,-15$, by using $j_t=1,2,2,2,3$.
Thus the second parenthesis is always at least $-20$.
Finally, for every integer $k$, with $x=k-28$,
\[
 q+3I=T_k+5k+3(k-34)^2
 =654+\frac{7x^2-5x}{2}\ge654.
\]
We have proved $L\ge492-3T$. By \eqref{eq:85-cross-fringe},
\[
 L+2(20-B)\ge492-3T+2\max(0,T-163)>0
 \qquad(T\le165).
\]
For $T\le163$ the margin is at least $3$; at $T=164,165$ it is at
least $2,1$, respectively. This contradicts \eqref{eq:85-cross-raw}
and excludes the conditional A0 cross case with first gap $1$.
